\chapter{Writing a Thesis}
\label{cha:TheThesis}

Every thesis%
\footnote{Most of the following remarks apply equally to bachelor's, master's,
and diploma theses.}
is different, yet good theses are usually very similar in structure, especially
in the fields of engineering and natural sciences.

\section{Elements of a Thesis}

The following basic structure has proven itself as a starting point, which can,
of course, be varied and refined as desired:
%
\begin{enumerate}
	\item \textbf{Introduction and motivation}: What is the problem statement or
	task at hand, and why should someone be interested in it?
	\item \textbf{Specification of the topic in greater detail}: Here, the
	current state of the art of the technology or science is described, and
	existing deficits or open questions are pointed out. The direction of one's
	work is developed from this.
	\item \textbf{Own approach}: This is, of course, the core of the thesis.
	Here it is shown how the previously described task is solved and
	implemented, often in the form of a prototype. Illustrative examples
	supplement this part.
	\item \textbf{Summary}: What has been achieved, and what goals remain open?
	Which parts of the thesis are possible origins for further work?
\end{enumerate}
%
Of course, a certain dramaturgical structure of the thesis is also important.
Remember that readers usually have little time and---unlike in a novel---their
patience should not be tried. Explain already in the introduction (not in the
last chapter) how you approached the problem, your proposed solutions, and if
you successfully applied them.

Errors and dead ends may (and should) be described as well; their knowledge
often helps to avoid duplicate experiments and other errors and is thus
certainly more helpful than any whitewashing. And, of course, it is acceptable
to express one's own ideas and opinions as long as they are rationally stated.


\section{Language and Writing Style}

A thesis is a piece of scientific work and should therefore be formulated concisely 
and factually. The author's person takes a back seat to the subject of the
work; avoid the first person or wordings such as "the author". Passive voice can
be a remedy, although it is vital to ensure that this does not result in an
overly complicated sentence structure. Where possible, prefer active wording
without the first person (\eg, "Chapter~3 describes \ldots" instead of "In
Chapter~3, I describe \ldots"), since it sounds more direct and delivers
essential aspects more precisely.

Expressions such as colloquialisms, polemical formulations, or even irony and
cynicism are out of place, as is the excessive use of overly specific technical
terms.

Furthermore, the language used in a thesis should be gender-inclusive and
non-discri\-minatory, making all people equally visible in their diversity,
both in words and images. In English, this includes using nouns that are not
gender-specific to refer to roles or professions, formation of phrases in a
coequal manner, and avoiding the blanket use of male or female terms.

Avoid gender-specific job titles such as "fireman" or "stewardess" and replace
them with the more neutral terms "firefighter" or "flight attendant". When
using pronouns, refrain from using "he" or "she". Replace these with the more
inclusive "they" to include people who identify as non-binary.

Also, check the writing for terms that might be considered racially
inappropriate. Expressions such as "master" and "slave" or "whitelist" and
"blacklist" might be considered offensive by certain groups of people. Keep this
in mind when naming things in the project or prototype; neutral alternatives
are, for example, "primary"/"replica" or "allowlist"/"blocklist". Especially in a
scientific thesis, the potential of language should be used to counter
stereotypical ideas about social roles.


\section{Writing a Thesis in English at a German-Speaking University}
\label{sec:german}

While this template and the contained introduction should make it easy to write
a thesis entirely in English, there may still be some things that have to be in
German, depending on the requirements of the university.

The title page is usually in the language of the thesis (in the past, a German
title page was mandatory). Its language can be set independently of the main
language with the option \texttt{titlelanguage} (see
Section~\ref{sec:hagenberg-settings}). Also, a German \emph{Kurzfassung} is
usually included together with the English \emph{Abstract}. If in doubt, check
both with the degree program. Using a translator is a valid option for everyone
who is not fluent in German. Having the translation checked by a native German
speaker is recommended, however.

The German term "Fachhochschule" (as in Fachhochschule Oberösterreich) is
translated with "University of Applied Sciences". A master's thesis is called
"Masterarbeit", and a bachelor's thesis is called "Bachelorarbeit". Academic
titles of persons are typically less important in English than in
German-speaking countries and are therefore usually omitted.

If one's native language is German, it should be considered that writing a
thesis in English (unless the program requires it) does not make writing any
easier, even if it might feel that way at first. Particularly in computer
science, the dominance of English technical terms makes writing in German seem
tedious, and switching to English might appear particularly attractive. However,
this is deceptive since one's skill in a foreign language is often overestimated
(despite the usually long years of English education). Conciseness and clarity
are easily lost, and sometimes the result is an embarrassing drivel without
context and solid content. Unless one's English skills are excellent, writing at
least the most essential parts of the thesis in German first and only
translating them afterward is advisable. Special care should be taken when
translating seemingly familiar technical terms. In addition, it is always
helpful to have the finished work checked by a native speaker.

Parallel to this document, there is also a German version, which is largely
identical in content. It contains helpful hints for writing a thesis in German
and should be used if German is the language of choice. Technically, except for
the language setting and the different quotation marks (see
Section~\ref{sec:quotation-marks}), there is nothing more to consider to use
this template in German.


\section{Making Use of AI Tools}

Artificial intelligence (AI) tools are now widely used for writing texts, this
includes academic writing. Many of these applications are based on so-called
\emph{large language models} (LLMs) that have been trained on huge amounts of text.
Through statistical analyses of millions or billions of words, these models
learn to generate meaningful sentences with high probability and to recognize
connections. Examples of AI-supported writing assistants include 
ChatGPT\footnote{\url{https://chatgpt.com/}},
DeepL Write\footnote{\url{https://www.deepl.com/write}} or 
Grammarly\footnote{\url{https://grammarly.com/}}.

When writing a thesis, AI tools can help make existing texts more efficient,
for example, by checking spelling and grammar, improving style and readability,
or suggesting initial wording. This saves time and allows the author to focus
more on the content quality and the work's structure.

At the same time, such tools also entail risks. In particular, they
\textbf{should not} be used to generate new content. There is a risk that
incorrect texts may be created or findings uncritically adopted. Furthermore,
AI models can reproduce prejudices or use lines of argument that contradict a
thorough scientific approach. Last but not least, the affidavit (as in the
documents created with this template) confirms that the thesis was "written
independently" and that "only the references and aids indicated" were used.
A text generated by an AI tool goes against the idea of a thesis as an
independent academic achievement.

\subsection{Important Points When Using AI Tools}

\begin{itemize}
	\item Use AI tools in the text only to revise the wording or to make
	linguistic improvements.
	\item State the use of AI tools in the documentation of AI tools usage at
	the beginning of the thesis (see Section~\ref{sec:DocumentingAIUsage}).
	\item If large parts of the text or revisions have been created with the
	help of such a tool, document this additionally at the respective
	position, for example, by adding a note in a footnote or at the beginning
	of a section. It is helpful to specify the prompt used if it goes beyond a
	simple "revise this paragraph".
	\item Adding a reference with the AI tool as the author should be avoided
	since the results are not deterministic.
	\item For AI-generated media (\eg, images) specify the tool used and the
	prompt in the caption. See Section \ref{sec:SourceCitationsInCaptions}.
	\item The responsibility for the AI content used always lies with the
	author. It is therefore important to critically review, question, and, if
	necessary, correct the results.
\end{itemize}

\subsection{Documenting the Use of AI Tools}
\label{sec:DocumentingAIUsage}

Many universities require the use of AI tools in a thesis to be documented. In
the documents created with this template, the affidavit is therefore followed
by a separate page \emph{Documentation of AI Tools Usage} with the text modules
of the University of Applied Sciences Upper Austria. It is created with the
environment \texttt{aiusage} in the file \verb!front/ai-usage.tex!, which is
included in \verb!main.tex! directly after \verb!\maketitle!. There, the
applicable text modules are selected and all others deleted, for example:
%
\begin{LaTeXCode}[numbers=none]
\begin{aiusage}
  \aiusedfor{literature-search}
  \aiusedfor{proofreading}
\end{aiusage}
\end{LaTeXCode}
%
If no AI tools were used, \verb!\aiusednone! is given instead. Further text
modules are:
%
\begin{itemize}
	\item \verb!\aiusedfor{ideas}!: gather initial ideas for research questions
	and theories,
	\item \verb!\aiusedfor{literature-understanding}!: better understand
	academic literature,
	\item \verb!\aiusedforother{...}!: other applications.
\end{itemize}
%
However, any types of use that do not appear on this list must be clarified
with your supervisor in advance.

If the degree program requires more detailed documentation, the long form
prepared in the commented part of \verb!front/ai-usage.tex! is used instead
(if both forms are used, \latex\ issues a warning).
It creates a form that confirms the responsible use
of AI tools and lists which AI tools were used for which activities (\eg,
literature search, coding, or proofreading):
%
\begin{LaTeXCode}[numbers=none]
\begin{aiusagedetails}
  \aiconfirm{accuracy}
  \aiconfirm{responsibility}
  \aiactivity{coding}{GitHub Copilot}{Generation of test functions (Chapter 4)}
  \aiactivity{editing}{DeepL Write}{Proofreading of Chapters 1--5}
  \aifurtheruse{Generation of test data for the evaluation (Section 5.2).}
\end{aiusagedetails}
\end{LaTeXCode}
%
Like the FH Upper Austria template, the table lists all activities by default;
with \verb!\begin{aiusagedetails}[compact]!, only those with entries are listed.
The instructions of the template (\eg, to name the AI tools) are not printed in
the form but given as comments in \verb!front/ai-usage.tex!. Uses that
correspond to none of the activities in the table are given with
\verb!\aifurtheruse{...}!, one command per use. A single entry is printed as a
paragraph, several entries as a bulleted list; without \verb!\aifurtheruse!,
"None" is printed.
All commands and activities are described in the package manual.
